% =====================================================================
% TOI-VELOCITY ADJOINT — HYBRID BACKWARD PASS (OPERATIVE TRACK, FILE 4)
% Adjoins: docs/toi_manifold.tex (surface + TOI gradient),
%          the formal envelope (M = <L,S,T>),
%          the Fibonacci/floating-base reference (Part II).
% Pure mechanics. No narrative constants. Compile: pdflatex toi_adjoint.tex
% =====================================================================
\documentclass[11pt]{article}
\usepackage{amsmath,amssymb,amsthm}
\usepackage[margin=1in]{geometry}
\usepackage{booktabs}

\newtheorem{proposition}{Proposition}
\newtheorem{definition}{Definition}

\title{The Adjoint of Time-of-Impact Hybrid Dynamics\\
      (Backward Pass for $\partial J/\partial u$)}
\date{2026-09-29}

\begin{document}
\maketitle

% ---------------------------------------------------------------------
\section{Problem}
% ---------------------------------------------------------------------
Minimise
\begin{equation}
  J(u) \;=\; \Phi(s(T)) \;+\; \int_0^T \! L(s,u)\,dt,
  \qquad s \in \mathbb{R}^{4d},
\end{equation}
subject to the hybrid dynamics of the adjoining manifold document:
free flow $\dot s = f(s,u)$ between impacts, and reset
$s^+ = g(s^-)$ at impact times $\gamma_k$ satisfying
$\psi(s(\gamma_k)) = 0$ with $\nabla\psi^\top f < 0$ (transversality).

% ---------------------------------------------------------------------
\section{The adjoint states}
% ---------------------------------------------------------------------
\begin{definition}[Continuous-phase adjoint]
Between impacts, the adjoint $(\lambda, \mu)$ obeys the standard
ODEs of the Pontryagin maximum principle:
\begin{align}
  \dot\lambda &= -\Bigl(\tfrac{\partial f}{\partial s}\Bigr)^\top \lambda
                 \;-\; \nabla_s L, \\
  \dot\mu     &= -\Bigl(\tfrac{\partial f}{\partial u}\Bigr)^\top \lambda
                 \;-\; \nabla_u L,
\end{align}
with terminal conditions $\lambda(T) = \nabla\Phi(s(T))$, $\mu(T) = 0$,
and the reduced gradient of $J$ with respect to the control on any
interval $(\gamma_k, \gamma_{k+1})$ is
$\partial J/\partial u = \mu$ evaluated on that interval.
\end{definition}

% ---------------------------------------------------------------------
\section{The jump map of the adjoint (the load-bearing step)}
% ---------------------------------------------------------------------
Naive hybrid adjoints carry only the Jacobian factor
$(\partial g/\partial s)^\top$ across the impact. That is wrong
whenever $\partial\gamma_k/\partial u \neq 0$, exactly as the forward
gradient is wrong without the transport term.

\begin{proposition}[Adjoint jump at an impact]
\label{prop:jump}
Integrating the adjoint backward across $\gamma_k$ gives
\begin{equation}
  \lambda(\gamma_k^-) \;=\;
  \Bigl(\tfrac{\partial g}{\partial s}\Bigr)^\top \lambda(\gamma_k^+)
  \;+\;
  \beta_k \, \nabla\psi\bigl(s(\gamma_k)\bigr),
\end{equation}
where the scalar impact multiplier $\beta_k$ is fixed by the
Hamiltonian jump condition at the impact (the adjoint image of
the forward transport term), and the control-side jump is
\begin{equation}
  \mu(\gamma_k^-) \;=\;
  \mu(\gamma_k^+) \;+\;
  \beta_k\,
  \frac{\nabla\psi^\top \tfrac{\partial s}{\partial u}}
       {\nabla\psi^\top f}\Big|_{\gamma_k}.
\end{equation}
\end{proposition}

Two remarks make the structure visible:

\begin{itemize}
  \item \textbf{The $\beta_k\,\nabla\psi$ term is the adjoint image of
        the forward transport term.} Forward:
        $f \otimes \partial\gamma_k/\partial u$ moves the impact time;
        backward: $\beta_k\,\nabla\psi$ moves the impact
        \emph{constraint}. They arise from the same differentiated
        manifold condition $\psi(s(\gamma_k))=0$.
  \item \textbf{Dropping $\beta_k$ is the adjoint false witness.}
        A backward pass that omits the multiplier reproduces exactly
        the published failure: pre-collision control that decreases
        when it should increase.
\end{itemize}

% ---------------------------------------------------------------------
\section{Discrete form (matches the code layer)}
% ---------------------------------------------------------------------
For the discretised problem with steps $\Delta t$ and a step map
that may contain one impact, the discrete adjoint of the
TOI-velocity step is the exact transpose of the forward
sensitivity — one multiplication by
$\bigl(\partial \operatorname{step}/\partial s\bigr)^\top$
per step, where the step Jacobian \emph{already contains} the
transport term of the adjoining document:
\begin{equation}
  \frac{\partial s_{i+1}}{\partial u_i}
  \;=\;
  \frac{\partial \operatorname{step}}{\partial s}\,
  \frac{\partial s_i}{\partial u_i}
  \;+\;
  \frac{\partial \operatorname{step}}{\partial \gamma}\,
  \frac{\partial \gamma}{\partial u_i},
  \qquad
  \frac{\partial \gamma}{\partial u_i}
  \;=\;
  -\,\frac{\nabla\psi^\top \tfrac{\partial s_i}{\partial u_i}}
          {\nabla\psi^\top f}.
\end{equation}
Backward-mode differentiation of this same expression
(autograd over the step) yields the adjoint automatically and
exactly; a correct TOI forward gradient and a correct TOI
adjoint are one object, not two.

% ---------------------------------------------------------------------
\section{Verification protocol (gate for the benchmark)}
% ---------------------------------------------------------------------
\begin{center}
\begin{tabular}{@{}lll@{}}
\toprule
Gate & Check & Assertion \\
\midrule
A1 & directional derivative & $\nabla J^\top \delta u \approx
      \bigl(J(u+\varepsilon\delta u) - J(u)\bigr)/\varepsilon$
      to $O(\varepsilon)$ \\
A2 & adjoint vs forward & adjoint gradient $\equiv$ forward-mode
      gradient (same TOI terms) \\
A3 & impact-count robustness & gates hold for $1..m$ impacts in
      $(0,T)$ \\
A4 & pre-collision $u_y$ & monotonically increasing (the
      discriminating signature, gate G2 of the benchmark) \\
\bottomrule
\end{tabular}
\end{center}
Gate A1 is the brute-force truth: finite differences on the
loss. Gate A2 catches transpose errors. Gate A3 catches
accumulated drift over multiple impacts. Gate A4 is inherited
from the adjoining benchmark and is the one naive adjoints fail.

% ---------------------------------------------------------------------
\section{Floating-base extension}
% ---------------------------------------------------------------------
For the legged robot the same jump applies at each touch-down,
with the manifold condition supplied by the unilateral foot
constraint and $g$ by the contact-impulse solution of
\begin{equation}
  M(q)\bigl(\ddot q + [0;g;0]\bigr) + n(q,\dot q)
  = u + \sum_i C_i(q)^\top f_i ,
\end{equation}
so the adjoint multiplier $\beta_k$ attaches to the foot-contact
constraint gradient, exactly as Proposition~\ref{prop:jump}
prescribes for balls. Every footstep is a TOI event; the adjoint
jump is the per-footstep correction.

\section{Ledger integration}
This document seals as \texttt{proof\_class: code} under
\texttt{GARDEN.EVENT.v1}. Gates A1--A4 are machine-checkable;
the workload \texttt{toivelocity-gradient-check} asserts A1--A4
alongside G1--G3 of the benchmark. All quantities are ordinary
rigid-body mechanics.

\end{document}
